\section{数值检验}
\label{sec:tests}

\begin{figure}[tb]
\includegraphics[width=0.78\textwidth]{images/stout_test_B57}
\caption[figA]{式~(\ref{eq:Eeff}) 定义的有效能量 $aE_{\rm eff}(r)$：
间距 $r=5a$ 的静态夸克-反夸克对在不同涂抹层数 $n_\rho=1,5,$ 与 $20$ 下
随涂抹参数 $\rho$ 的变化。
 结果在 $12^4$ 格点上用耦合 $\beta=5.7$ 的 Wilson 规范作用量得到。
 标记 $A_{n_\rho}$ 的曲线是采用解析 stout 链接涂抹方案的空间各向同性
 三维版本（$n_\rho=1,5,$ 与 $20$ 层）的结果，而标记 $B_{n_\rho}$ 的曲线
 是用式~(\ref{eq:Ufuzz}) 配合文献~\cite{su3projectA} 的投影法涂抹链接的
 结果。
\label{fig:stout_veff}}
\end{figure}

我们在若干蒙特卡罗模拟中检验了 stout 链接的双重功效：既作为胶子算符的
涂抹算法，也用于在短距离量中压低紫外胶子模式的效应。

间距为 $r$ 的静态夸克-反夸克对存在时的胶子能量，可以在时间延伸 $t$
变大时从 $r\times t$ Wilson 圈 $W(r,t)$ 中提取。Wilson 圈可视为一个规范
不变算符的关联函数，该算符由一对静态夸克-反夸克以及沿夸克与反夸克之间
直线路径排列的链接变量乘积构成。若该夸克-反夸克-胶子算符与较高能级的
耦合很小，最低能量的提取会可靠得多。这可以通过用连接夸克与反夸克的涂抹
链接乘积代替原始链接变量来实现。换言之，静态夸克-反夸克对存在时胶子的
基态能量可以从 Wilson 圈 $\tilde{W}(r,t)$ 更可靠地确定，其中初始与末端
时间片上 $r$ 链接路径使用的是涂抹的空间链接。衡量到较高能级的耦合被压低
程度的量度，是单时间步的有效能量，定义为
\begin{equation}
  a E_{\rm eff}(r) = -\ln\left(\frac{\tilde{W}(r,a)}{\tilde{W}(r,0)}
 \right).
 \label{eq:Eeff}
\end{equation}
对具有正定转移矩阵的格点规范作用量（例如 Wilson 作用量），在通常情况下，
减少激发态污染往往使上述有效能量降低。


\begin{figure}[tb]
\includegraphics[width=0.78\textwidth]{images/stout_test_B62}
\caption[figB]{式~(\ref{eq:Eeff}) 定义的有效能量 $aE_{\rm eff}(r)$：
间距 $r=10a$ 的静态夸克-反夸克对在不同涂抹层数 $n_\rho=1,5,$ 与 $20$ 下
随涂抹参数 $\rho$ 的变化。
 结果在 $24^4$ 格点上用耦合 $\beta=6.2$ 的 Wilson 规范作用量得到。
 标记 $A_{n_\rho}$ 的曲线是采用解析 stout 链接涂抹方案的空间各向同性
 三维版本（$n_\rho=1,5,$ 与 $20$ 层）的结果，而标记 $B_{n_\rho}$ 的曲线
 是用式~(\ref{eq:Ufuzz}) 配合文献~\cite{su3projectA} 的投影法涂抹链接的
 结果。
\label{fig:stout_veff2}}
\end{figure}

我们用式~(\ref{eq:Eeff}) 定义的有效能量检验了 stout 链接涂抹方案抑制
胶子算符中激发态污染的能力。$12^4$ 格点、$\beta=5.7$、$r=5a$ 的结果示于
图~\ref{fig:stout_veff}；更大的 $24^4$ 格点、$\beta=6.2$、$r=10a$ 的结果
示于图~\ref{fig:stout_veff2}。
这些结果与用式~(\ref{eq:Ufuzz}) 及文献~\cite{su3projectA} 的投影法涂抹
链接所得的结果作了比较。

两幅图都表明，链接涂抹能够显著减少来自理论高激发模式的污染。先看采用
stout 链接涂抹方案的空间各向同性三维版本（$\rho_{jk}=\rho,
\ \rho_{4\mu}=\rho_{\mu 4}=0$）得到的结果。 对给定的
涂抹迭代次数 $n_\rho$，当涂抹参数 $\rho$ 从零开始增大时，有效能量起初
下降。 最终会达到某个最优值，在该处有效能量取极小。
这一最优值随涂抹层数 $n_\rho$ 增大而减小，且在这个最优 $\rho$ 处有效
能量的下降幅度相当可观。 继续把 $\rho$ 增大到超过其
最优值后，有效能量急剧上升。
围绕极小值的下降与上升的陡峭程度，对更大的 $n_\rho$ 更为显著。此外，
极小值随 $n_\rho$ 增大而降低，直至达到饱和点，超过该点后不再有额外的
压低。

用式~(\ref{eq:Ufuzz}) 加文献~\cite{su3projectA} 投影法涂抹链接所得到的
结果显示出基本相同的趋势，只是有效能量的每个极小值要宽得多。 stout
链接涂抹方案对参数 $\rho$ 的敏感性增大并不奇怪，因为 $\rho$ 出现在指数
函数内部。 然而，值得注意的是，两种涂抹方法给出的有效能量极小值几乎
相同。 就抑制胶子算符中的激发态污染而言，stout 链接涂抹方案与现行的
标准链接 fuzz 方案同样有效，只是由于对涂抹参数 $\rho$ 的敏感性增大而需要
更仔细的调校。
注意，这种解析链接涂抹方法已成功应用于 torelon 激发谱的计算%
\cite{torelons}。

\begin{figure}[tb]
\includegraphics[width=0.78\textwidth]{images/stout4D_plaq_B57}
\caption[figC]{平均涂抹方格值随涂抹参数 $\rho$ 的变化。结果用 Wilson 规范
作用量在 $12^4$ 格点上、耦合 $\beta=5.7$ 时得到。标记 $A_{n_\rho}$ 的曲线
是解析 stout 链接涂抹方案的各向同性四维版本（$n_\rho=1$ 与 $5$ 层）的
结果，而标记 $B_{n_\rho}$ 的曲线是用式~(\ref{eq:Ufuzz}) 配合文献%
\cite{su3projectA} 的投影法涂抹链接的结果。
\label{fig:stout_plaq}}
\end{figure}

在格点规范与费米子作用量中使用涂抹链接的主要目的，是减小格点理论中紫外
模式引起的离散效应。这类效应常被蝌蚪图的大贡献主导。 衡量一个链接
涂抹方案能在多大程度上减少蝌蚪图假效应的一个简单量度，是平均涂抹方格值。
平均涂抹方格值接近 1 表明蝌蚪图贡献被大幅削减。

$12^4$ 格点、Wilson 规范作用量、耦合 $\beta=5.7$ 的平均涂抹方格值结果
示于图~\ref{fig:stout_plaq}。 对两种涂抹方法的各向同性四维版本
（$\rho_{\mu\nu}=\rho$），当 $\rho$ 从零开始增大时，平均涂抹方格值起初
向 1 上升，直至达到极大值。 两种方法给出的极大值几乎相同。 $\rho$
继续增大时，stout 链接方案的平均方格值很快开始回落，而标准方案几乎不变。
$n_\rho$ 从零开始增大时，平均方格值的极大值也随之升高，但最终达到
饱和点。 注意 $n_\rho=5$ 的极大值非常接近 1，表明蝌蚪图贡献被急剧
压低。
总之，可以观察到：解析 stout 链接涂抹方案在减少紫外胶子模式的离散效应
方面与标准涂抹方案同样有效，但对涂抹参数 $\rho$ 的敏感性增大。
